
Page 13 resist. compression and distortion, he refused to accept the conclusion
from Poisson's theory that the modulus of compression has to the rigidity the
ratio \(5:3\). He pointed out that, if the ratio of these moduluses
could he regarded as infinite, the ratio of the velocities of "longitudinal"
and "transverse" waves would also be infinite, and then, as Green had already
shown, the application of the theory to optics would be facilitated.


The methods of Navier, of Poisson, and of Cauchy's later memoirs lead to
equations of motion containing fewer constants than occur in the equations
obtained by the methods of Green, of Stokes, and of Cauchy's first memoir. The
importance of the discrepancy was first emphasized by Stokes. The questions in
dispute are these — Is elastic aeolotropy to be characterized by 21 constants
or by 15, and is elastic isotropy to be characterized by two constants or one ?
The two theories are styled by Pearson the "multi-constant" theory and the
"rari-constant" theory respectively, and the controversy concerning them has
lasted almost down to the present time. It is to be understood that the
rari-constant equations can be included in the multi-constant ones by equating
certain pairs of the coefficients, but that the rariconstant equations rest
upon a particular hypothesis concerning the constitution of matter, while the
adoption of multi-constancy has been held to imply denial of this hypothesis.
Discrepancies between the results of the two theories can be submitted to the
test of experiment, and it might be thought that the verdict would be final,
but the difficulty of being certain that the tested material is isotropic has
diminished the credit of many experimental investigations, and the tendency of
the multi-constant elasticians to rely on experiments on such bodies as cork,
jelly and india-rubber has weakened their arguments, much of the discussion
has turned upon the value of the ratio of lateral contraction to longitudinal
extension of a bar under terminal tractive load. This ratio is often called
"Poisson's ratio."  Poisson \footnote{Toihniitera.nAVearsoD,HistoryoftheTheoryof
Elasticity, Yo. 1, Cambridge, 1886, p. 496.}> deduced from his theory the
result that this ratio must be \(1/4\). The experiments of Wertheim
\footnote{ Annates de Chimie, t. 23 (1848)}> on glass and brass did not
support this result, and Wertheim proposed to take the ratio to be \(1/3\) —
a value which has no theoretical foundation. The
experimental evidence led Lam in his treatise\footnote{Legons sur la thiorie
mathematique de Velasticite des corps solides, Paris, 1852.}> to adopt the
multi-constant equations, and after the publication of this book they were
generally employed. Saint-Venant, though a firm believer in rari-constancy,

Page 14 of the atoms, we can only invoke the known laws of energy as was done by
Green and Lord Kelvin; and we may place the theory on a firm basis if we
appeal to experiment to support the statement that, within a certain range
of strain, the strain-energy-function is a quadratic function of the
components of strain, instead of relying, as Green did, upon an expansion
of the function in series.

The problem of determining the state of stress and strain within a solid
body which is subjected to given forces acting through its volume and to
given tractions across its surface, or is held by surface tractions so that its
surface is deformed into a prescribed figure, is reducible to the analytical
problem of finding functions to represent the components of displacement.
These functions must satisfy the differential equations of equilibrium at all
points within the surface of the body and must also satisfy certain special
conditions at this surface. The methods which have been devised for integrating
the equations fall into two classes. In one class of methods a special
solution is sought and the boundary conditions are satisfied by a solution in
the form of a series, which may be infinite, of special solutions. The special
solutions are generally expressible in terms of harmonic functions. This class
of solutions may be regarded as constituting an extension of the methods of
expansion in spherical harmonics and in trigonometrical series. In the other
class of methods the quantities to be determined are expressed by definite
integrals, the elements of the integrals representing the effects of
\texttt{singularities}
distributed over the surface or through the volume. This class of solutions
constitutes an extension of the methods introduced by Green in the Theory
of the Potential. At the time of the discovery of the general equations of
Elasticity the method of series had already been applied to astronomical
problems, to acoustical problems and to problems of the conduction of
heat\footnote{See Burkhardt, `Ent wickelungen nach oscillirenden Functionen'
Jahresbericht der Deutschen Mathematiker-Vereinigung, Bd. 10, Heft 2.};
the method of singularities had not been
invented\footnote{It was invented by Green, An Essay on the Application of
Mathematical Analysis to the Theories of Electricity and Magnetism, Nottingham,
1828. Reprinted in Mathematical Papers of the late George Green, London,
1871.}. The application of the method of series to problems of equilibrium of
elastic solid bodies was initiated
\setcounter{footnote}{38}
by Lamé and Clapeyron\footnote{See footnote 39}. They considered the case of a
body bounded by an unlimited plane to which pressure is applied according to an
arbitrary law.
\setcounter{footnote}{59}
Lamé\footnote{Journal de Math. (Liouville), t. 19 (1854).}
later considered the problem of a body bounded by a spherical surface
and deformed by given surface tractions. The problem of the plane is
essentially that of the transmission into a solid body of force applied locally
to a small part of its surface. The problem of the sphere has been developed
by Lord
Kelvin\footnote{Phil. Trans. Boy. Soc, vol. 153 (1863). See also Mathematics
and Phys. Papers, vol. 3 (Cambridge, 1890), p. 351, and Kelvin and Tait, Nat.
Phil., Part 2.} who sought to utilize it for the purpose of investigating the
rigidity of the Earth\footnote{Brit. Assoc. Rep. 1876, Math, and Phys. Papers,
vol. 3, p. 312.}, and by G. H. Darwin in connexion with other problems of
cosmical
physics\footnote{Phil. Trans. Roy. Soc, vol. 170 (1879), and vol. 173 (1882).
Reprinted in G. H. Darwin's Scientific Papers, vol. 2, Cambridge 1908, pp. 1, 459.}.
The serial solutions employed are expressed
in terms of spherical harmonics. Solutions of the equations in cylindrical
coordinates can be expressed in terms of Bessel's
functions\footnote{L. Pochhammer, J.f. Math. (Crelle), Bd. 81 (1876), p. 38.},
but, except for spheres and cylinders, the method of series has not been employed very
successfully. The method of singularities was first applied to the theory of
Elasticity by E. Betti\footnote{II Nuovo Cimento (Ser. 2), tt. 6—10 (1872 et seq.)},
who set out from a certain reciprocal theorem of the
type that is now familiar in many branches of mathematical physics. From
this theorem he deduced incidentally a formula for determining the average
strain of any type that is produced in a body by given forces. The method
of singularities has been developed chiefly by the elasticians of the Italian
school. It has proved more effective than the method of series in the solution
of the problem of transmission of force. The fundamental particular solution
which expresses the displacement due to force at a point in an indefinitely
extended solid was given by Lord
Kelvin\footnote{Sir W. Thomson, Cambridge and Dublin Math. J., 1848, reprinted in Math
and Phys. Papers, vol. 1, p. 97}.  It was found at a later date by J.
Boussinesq\footnote{For Boussinesq's earlier researches in
regard to simple solutions, see Paris, C. R., tt. 86—88 (1878—1879)
and tt. 93—96 (1881—1883). A more complete account is given in his book,
Applications des potentiels à l'étude de l'équilibre et du mouvement
des solides elastiques, Paris, 1885.}
along with other particular solutions, which can, as a matter
of fact, be derived by synthesis from it. Boussinesq's results
led him to a solution of the problem of the plane, and to a theory of
``local perturbations'' according to which the effect of force applied
in the neighbourhood of any point of a body falls off very rapidly as
the distance from the point increases, and the application of an
equilibrating system of forces to a small part of a body produces an
effect which is negligible at a considerable distance from the part.
To estimate the effect produced at a distance by forces applied
near a point, it is not necessary to take into account the mode of
application of the forces but only the statical resultant and moment.
The direct method of integration founded upon Betti's reciprocal
theorem was applied to the problem of the plane by V.
Cerruti\footnote{Rome, Acc. Lincei, Mem.\ fis.\ mat., 1882.
}. Some of the results were found independently by Hertz, and
led in his hands to a theory of impact and a theory of
hardness\footnote{J. f. Math. (Crelle), Bd. 92 (1882), and Verhandlungen
des Vereins zur Beförderung des Gewerbefleisses, Berlin, 1882. The memoirs
are reprinted in Ges. Werke von Heinrich Hertz, Bd. 1, Leipzig, 1895, pp. 155 and 174.
}.

A different method for determining the state of stress in a body
has been developed from a result noted by G. B.
Airy\footnote{70 Brit. Assoc. Rep. 1862, and Phil. Trans. Roy. Soc, vol. 153
(1863), p. 49}.  He observed that, in the case
of two dimensions, the equations of equilibrium of a body deformed by surface
tractions show that the stress-components can be expressed as partial differential
coefficients of the second order of a single function.
Clerk Maxwell\footnote{71 Edinburgh Roy. Soc. Trans., vol. 26 (1870). Reprinted in Maxwell's
Scientific Papers, vol. 2, p. 161.} ``extended the
result to three dimensions, in which case three such ``stress-functions'' are
required. It appeared later that these functions are connected by a rather
complicated system of, differential equations\footnote{72 W. J. Ibbetson, An Elementary
Treatise on the Mathematical Theory of
perfectly Elastic Solids, London, 1887.}. The stress-components must
in fact be connected with the strain-components by the stress-strain relations,
and the strain-components are not independent; but the second differential
coefficients of the strain-components with respect to the coordinates are
connected by a system of linear equations, which are the conditions necessary to
secure that the strain-components shall correspond with a displacement, in
accordance with the ordinary formulas connecting strain and
displacement\footnote{73 Saint-Venant gave the identical relations between
strain-components in his edition of Navier's Résumé des Leçons sur I'application
de la Mécanique, Paris, 1864, Appendice 3}.
It is possible by taking account of these relations to obtain a complete system
of equations which must be satisfied by stress-components, and thus the way
is open for a direct determination of stress without the intermediate steps of
forming and solving differential equations to determine the components of
displacement\footnote{J. H. Michell, London Math. Soc. Proc, vol. 31 (1900), p. 100.}.
In the case of two dimensions the resulting equations are of a
simple character, and many interesting solutions can be obtained.

The theory of the free vibrations of solid bodies requires the integration
of the equations of vibratory motion in accordance with prescribed boundary
conditions of stress or displacement. Poisson\footnote{36} gave the solution of the
problem of free radial vibrations of a solid sphere, and Clebsch\footnote{53}
founded the general theory on the model of Poisson's solution. This theory
included the extension of the notion of ``principal coordinates'' to systems
with an infinite number of degrees of freedom, the introduction of the
corresponding ``normal functions'' and the proof of those properties of these
functions upon which the expansions of arbitrary functions depend. The
discussions which had taken place before and during the time of Poisson
concerning the vibrations of strings, bars, membranes and plates had prepared
the way for Clebsch's generalizations. Before the publication of Clebsch's
treatise a different theory had been propounded by Lamé\footnote{49}.
Acquainted with Poisson's discovery of two types of waves, he concluded that
the vibrations of any solid-body must fall into two corresponding classes, and
he investigated the vibrations of various bodies on this assumption. The fact
that his solutions do not satisfy the conditions which hold at the boundaries
of bodies free from surface traction is a sufficient disproof of his theory;
but it was finally disposed of when all the modes of free vibration of a
homogeneous isotropic sphere were determined, \setcounter{footnote}{74}
and it was proved that the classes into which they fall do not verify
Lamé's supposition. The analysis of the general problem of the vibrations
of a sphere was first completely given by 
P. Jaerisch\footnote{75 J.f. Math. (Crelle), Bd. 88 (1880)},
who showed that the
solution could be expressed by means of spherical harmonics and certain
functions of the distance from the centre of the sphere, which are practically
Bessel's functions of order integer \(+1/2\). This result was obtained
independently by H. Lamb\footnote{76 London Math. Soc. Proc, vol. 13 (1882)},
who gave an account of the simpler modes of vibration
and of the nature of the nodal division of the sphere which occurs when any
normal vibration is executed. He also calculated the more important roots
of the frequency equation. 
L. Pochhammer\footnote{77 J. f. Math. (Crelle), Bd. 81 (1876), p. 324
}has applied the method of normal
functions to the vibrations of cylinders, and has found modes of vibration
analogous to the known types of vibration of bars.

The problem of tracing, by means of the equations of vibratory motion,
the propagation of waves through an elastic solid medium requires investigations
of a different character from those concerned with normal modes of
vibration. In the case of an isotropic medium
Poisson\footnote{78 Paris, Mém.\ de l'Acad.,t. 10 (1831).}
and Ostrogradsky\footnote{79 St Petersburg, Mém.\ de I' Acad., t. 1 (1831).} adopted
methods which involve a synthesis of solutions of simple harmonic
type, and obtained a solution expressing the displacement at any time in
terms of the initial distribution of displacement and velocity. The investigation
was afterwards conducted in a different fashion by
Stokes\footnote{80. `On the Dynamical Theory of Diffraction'
Cambridge Phil. Soc. Trans., vol. 9 (1849). Reprinted in Stokes's Math, and Phys.
Papers, vol. 2 (Cambridge, 1883)}, who showed
that Poisson's two waves are waves of irrotational dilatation and waves of
equivoluminal distortion, the latter involving rotation of the elements of the
medium. Cauchy\footnote{41} and
\setcounter{footnote}{80}
Green\footnote{81 Cambridge Phil. Soc. Trans., vol. 7 (1839).  Reprinted in
Green's Mathematical Papers, p. 293} discussed the propagation of plane waves
through a crystalline medium, and obtained equations for the velocity of
propagation in terms of the direction of the normal to the wave-front. In general
the wave-surface has three sheets; when the medium is isotropic all the sheets
are spheres, and two of them are coincident. Blanchet\footnote{82 82 J. de Math.
(Liouville), t. 5 (1840), t. 7 (1842).
}
extended Poisson's
results to the case of a crystalline medium. Christoffel\footnote{83 83 Ann.\ di Mat.
(Ser. 2), t. 8 (1877).Reprinted in E. B. Christoffel Ges.\ math. Abhandlmgen, Bd. 2,
p. 81, Leipzig 1910.} discussed the advance
through the medium of a surface of discontinuity. At any instant,
the surface separates two portions of the medium in which the displacements are
expressed by different formulae; and Christoffel showed that the surface moves
normally to itself with a velocity which is determined, at any point, by the
direction of the normal to the surface, according to the same law as holds

